\documentclass[11pt]{article}
\usepackage[T1]{fontenc}
\usepackage{lmodern}
\usepackage[margin=1in]{geometry}
\usepackage{amsmath,amssymb,amsthm,mathtools}
\usepackage[hidelinks]{hyperref}
\usepackage{microtype}
\newtheorem{theorem}{Theorem}
\newtheorem{lemma}{Lemma}
\newcommand{\R}{\mathbb R}
\newcommand{\norm}[1]{\lVert#1\rVert}
\setlength{\emergencystretch}{3em}
\title{Endpoint sharpness in two Grand Lebesgue inequalities\\
via concentration and Gaussian dilations}
\author{Candidate partial result from run \texttt{fa\_banach\_001}, lane 18}
\date{11 August 2026}

\begin{document}
\maketitle
\begin{abstract}
Liflyand, Ostrovsky, and Sirota left open the sharpness of a Sobolev
trace estimate and a Young convolution estimate in bilateral Grand Lebesgue
spaces. Their fixed borderline examples lose respectively $1/n$ and one unit
in the endpoint blow-up exponent. We show that these losses are artifacts of
using one fixed input. Normalized smooth concentration families and normalized
Gaussian families attain, up to constants, the complete endpoint orders
predicted by the two source inequalities for two-sided power gauges. The same
argument permits positive-index slowly varying endpoint modifiers. Exact
constants and completely arbitrary generating functions are not addressed.
\end{abstract}

\section{Source question and scope}
For a positive generating function $\psi$ on $(a,b)$, write
\[
 \norm{f}_{G(\psi;a,b)}
 :=\sup_{a<p<b}\frac{\norm{f}_{L^p}}{\psi(p)}.
\]
The source paper arXiv:0801.3404 proves a Sobolev trace estimate, labelled
(in2), and a convolution estimate, labelled (in3). In its final Sharpness
section it says (source PDF p.~11):
\begin{quote}
\emph{Finding sharp estimates in these two cases is an interesting open problem.}
\end{quote}
For the power gauge
\[
 \psi_{a,b}^{\alpha,\beta}(p)
 :=(p-a)^{-\alpha}(b-p)^{-\beta},\qquad a<p<b,
 \tag{1}
\]
the source's fixed examples give the correct endpoint interval but fall short
of the upper estimate by $1/n$ in (in2) and by $1$ in (in3).

The result below closes both exponent gaps in the natural operator-norm sense:
the test function may depend on the exponent tending to the endpoint. We work
with the standard restriction trace on $\R^m\subset\R^n$ and with convolution
on $\R^d$. These geometries already prove that the source estimates cannot be
improved in general.

\section{Proof intuition}
A fixed borderline singularity mixes growth of its input norm with the
smoothing caused by restriction or convolution. It therefore need not measure
the operator norm near every endpoint. Instead, move a smooth profile through
its dilation orbit.

For the trace inequality, $u_t(x)=c_t\phi(tx)$ has
\[
 \norm{\nabla u_t}_p=c_t t^{1-n/p}\norm{\nabla\phi}_p,\qquad
 \norm{u_t|_{\R^m}}_q=c_t t^{-m/q}\norm{\phi|_{\R^m}}_q.
\]
Normalizing the first expression in the Grand Lebesgue norm introduces
$(\log t)^\beta$ as $t\to\infty$. The remaining power of $t$ is
$\exp[-m(1/q-1/B)\log t]$. Choosing
$\log t\asymp(B-q)^{-1}$ recovers the missing endpoint power.

For convolution, Gaussians are closed under dilation and convolution. Their
normalizers contribute $(\log t)^{\beta_1+\beta_2}$, while their convolved
$L^r$ norm contributes
$\exp[-(d/2)(1/r-1/B)\log t]$. The same optimization recovers the whole
predicted exponent. Broad Gaussians give the lower endpoints.

\section{A normalizer lemma}
We use $F\asymp G$ for two-sided estimates with constants independent of the
parameter tending to the indicated endpoint.

\begin{lemma}[endpoint Laplace principle]\label{lem:laplace}
Let $0<a<b<\infty$, $\alpha,\beta>0$, and let $c$ be positive and
continuous on $[a,b]$. Put
\[
 M_\rho(t)=\sup_{a<p<b}
 c(p)t^{\rho-\kappa/p}(p-a)^\alpha(b-p)^\beta,\qquad \kappa>0.
\]
Then
\begin{align*}
 M_\rho(t)&\asymp
 t^{\rho-\kappa/b}(\log t)^{-\beta},&&t\to\infty,\\
 M_\rho(t)&\asymp
 t^{\rho-\kappa/a}(\log(1/t))^{-\alpha},&&t\downarrow0.
\end{align*}
\end{lemma}

\begin{proof}
For $t=e^L\to\infty$, set $s=b-p$. After factoring out
$t^{\rho-\kappa/b}$, the remaining exponential is
\[
 \exp\!\left[-\kappa L\left(\frac1p-\frac1b\right)\right]
 =\exp\!\left[-\frac{\kappa Ls}{bp}\right].
\]
On a fixed neighborhood of $b$, all other factors except $s^\beta$ are
bounded above and below. Thus the supremum is comparable to
$\sup_{0<s<s_0}s^\beta e^{-cLs}\asymp L^{-\beta}$. Outside that
neighborhood it is exponentially smaller. The proof at zero is identical,
with $s=p-a$ and $L=\log(1/t)$.
\end{proof}

\section{Sharp Sobolev trace endpoints}
Let $T u=u|_{\R^m}$ be restriction to
$\R^m\times\{0\}\subset\R^n$. Fix
\[
 1\le a<b<n,\qquad
 A=\frac{ma}{n-a},\qquad B=\frac{mb}{n-b},
 \tag{2}
\]
and assume $1\le A<B$. Let $p(q)=qn/(q+m)$, so $p(A)=a$,
$p(B)=b$, and define the source target gauge
\[
 \nu(q)=q^{1-1/n}\psi_{a,b}^{\alpha,\beta}(p(q)).
 \tag{3}
\]

\begin{theorem}[both trace endpoint exponents are sharp]\label{thm:trace}
For
\[
 \mathcal S(q):=\sup\left\{\norm{Tu}_{L^q(\R^m)}:
 u\in C_c^\infty(\R^n),\
 \norm{\nabla u}_{G(\psi_{a,b}^{\alpha,\beta};a,b)}\le1\right\},
\]
the source upper estimate and matching lower estimates give
\begin{align*}
 \mathcal S(q)&\asymp \nu(q)\asymp(q-A)^{-\alpha},
 &&q\downarrow A,\\
 \mathcal S(q)&\asymp \nu(q)\asymp(B-q)^{-\beta},
 &&q\uparrow B.
\end{align*}
Consequently neither endpoint power in (in2) can be decreased.
\end{theorem}

\begin{proof}
The upper bounds are (in2), and $p'(q)=nm/(q+m)^2>0$ shows that
$p(q)-a\asymp q-A$ and $b-p(q)\asymp B-q$.

Choose $\phi\in C_c^\infty(\R^n)$ with $T\phi\ne0$. Put
$u_t=c_t\phi(t\,\cdot)$ and choose $c_t$ so that
$\norm{\nabla u_t}_{G(\psi)}=1$. Since
\[
 \norm{\nabla(\phi(t\,\cdot))}_p
 =t^{1-n/p}\norm{\nabla\phi}_p,
\]
Lemma~\ref{lem:laplace}, with $\rho=1$ and $\kappa=n$, gives
\begin{align*}
 c_t&\asymp t^{n/b-1}(\log t)^\beta,&&t\to\infty,\\
 c_t&\asymp t^{n/a-1}(\log(1/t))^\alpha,&&t\downarrow0.
\end{align*}
Also
\[
 \norm{Tu_t}_q=c_t t^{-m/q}\norm{T\phi}_q,
\]
where $\norm{T\phi}_q$ stays bounded above and below on $[A,B]$.
The identities $n/b-1=m/B$ and $n/a-1=m/A$ yield
\begin{align*}
 \norm{Tu_t}_q&\asymp
 t^{-m(1/q-1/B)}(\log t)^\beta,&&t\to\infty,\\
 \norm{Tu_t}_q&\asymp
 t^{m(1/A-1/q)}(\log(1/t))^\alpha,&&t\downarrow0.
\end{align*}
For $q<B$, choose $\log t=(1/q-1/B)^{-1}$. For $q>A$, choose
$\log(1/t)=(1/A-1/q)^{-1}$. The reciprocal differences are comparable
to $(B-q)^{-1}$ and $(q-A)^{-1}$, proving the lower bounds.
\end{proof}

\section{Sharp convolution endpoints}
For $i=1,2$, let
\[
 \psi_i(p)=(p-a_i)^{-\alpha_i}(b_i-p)^{-\beta_i},
 \quad 1\le a_i<b_i<\infty,\quad \alpha_i,\beta_i>0.
\]
Assume
\[
 \frac1{a_1}+\frac1{a_2}>1,\qquad
 \frac1{b_1}+\frac1{b_2}>1,
\]
and define
\[
 \frac1A=\frac1{a_1}+\frac1{a_2}-1,\qquad
 \frac1B=\frac1{b_1}+\frac1{b_2}-1.
 \tag{4}
\]
For $A<r<B$, let
\[
 \tau(r)=\inf_{\substack{a_1<p<b_1,\ a_2<q<b_2\\
                    1/p+1/q=1+1/r}}
 \psi_1(p)\psi_2(q),
 \tag{5}
\]
the target gauge in (in3).

\begin{theorem}[both convolution endpoint exponents are sharp]\label{thm:conv}
Let
\[
 \mathcal C(r)=\sup\left\{\norm{f*g}_{L^r(\R^d)}:
 \norm f_{G(\psi_1)}\le1,\ \norm g_{G(\psi_2)}\le1\right\}.
\]
Then
\begin{align*}
 \mathcal C(r)&\asymp\tau(r)\asymp
 (r-A)^{-(\alpha_1+\alpha_2)},&&r\downarrow A,\\
 \mathcal C(r)&\asymp\tau(r)\asymp
 (B-r)^{-(\beta_1+\beta_2)},&&r\uparrow B.
\end{align*}
Thus the source's one-unit fixed-function gap does not occur for the
bilinear operator norm.
\end{theorem}

\begin{proof}
The upper bound $\mathcal C(r)\le\tau(r)$ is (in3). Near $B$, write
$p=b_1-s_1$ and $q=b_2-s_2$. The constraint (5) is exactly
\[
 \frac{s_1}{b_1p}+\frac{s_2}{b_2q}=\frac1r-\frac1B.
 \tag{6}
\]
Thus $s_1+s_2\asymp B-r$. Minimizing
$s_1^{-\beta_1}s_2^{-\beta_2}$ subject to (6) gives
$\tau(r)\asymp(B-r)^{-(\beta_1+\beta_2)}$. Near $A$, use
$p=a_1+s_1$, $q=a_2+s_2$, and $1/A-1/r\asymp r-A$.

For the lower bounds put $h_t(x)=e^{-t|x|^2}$ and
$N_i(t)=\norm{h_t}_{G(\psi_i)}$. Since
\[
 \norm{h_t}_p=\left(\frac{\pi}{pt}\right)^{d/(2p)},
\]
Lemma~\ref{lem:laplace} gives
\begin{align*}
 N_i(t)&\asymp t^{-d/(2b_i)}(\log t)^{-\beta_i},&&t\to\infty,\\
 N_i(t)&\asymp t^{-d/(2a_i)}(\log(1/t))^{-\alpha_i},&&t\downarrow0.
\end{align*}
The functions $f_t=h_t/N_1(t)$ and $g_t=h_t/N_2(t)$ have norm one.
The exact identity
\[
 h_t*h_t(x)=\left(\frac{\pi}{2t}\right)^{d/2}e^{-t|x|^2/2}
\]
implies
\begin{align*}
 \norm{f_t*g_t}_r
 &\asymp t^{-\frac d2(1/r-1/B)}
          (\log t)^{\beta_1+\beta_2},&&t\to\infty,\\
 \norm{f_t*g_t}_r
 &\asymp t^{\frac d2(1/A-1/r)}
          (\log(1/t))^{\alpha_1+\alpha_2},&&t\downarrow0.
\end{align*}
Choose $\log t=(1/r-1/B)^{-1}$ in the first line and
$\log(1/t)=(1/A-1/r)^{-1}$ in the second. This gives the asserted
lower orders.
\end{proof}

\section{Regularly varying upgrade}
The mechanism is not restricted to pure powers. Suppose near an upper
endpoint that
\[
 \psi(p)\asymp(b-p)^{-\beta}\ell\!\left(\frac1{b-p}\right),
 \qquad \beta>0,
 \tag{7}
\]
where $\ell$ is positive, locally bounded together with its reciprocal, and
slowly varying at infinity. Potter bounds in Lemma~\ref{lem:laplace} give an
extra factor $1/\ell(\log t)$ in the normalizer. Therefore
\[
 \mathcal S(q)\asymp
 (B-q)^{-\beta}\ell\!\left(\frac1{B-q}\right).
\]
For convolution, the two slowly varying factors multiply:
\[
 \mathcal C(r)\asymp\tau(r)\asymp
 (B-r)^{-(\beta_1+\beta_2)}
 \ell_{1,+}\!\left(\frac1{B-r}\right)
 \ell_{2,+}\!\left(\frac1{B-r}\right).
\]
The lower endpoints are analogous. Positive power index keeps the maximizing
dilation at distance comparable to $1/\log t$ from the endpoint; index-zero
gauges can require a different conjugate scale.

\section{Verification, novelty, and limitations}
The proof uses exact dilation identities, the elementary Gaussian convolution
formula, and the source upper estimates. Every exponent follows from
\[
 \frac nb-1=\frac mB,\quad
 \frac na-1=\frac mA,\quad
 \frac1B=\frac1{b_1}+\frac1{b_2}-1,\quad
 \frac1A=\frac1{a_1}+\frac1{a_2}-1.
\]
The optimizing scales tend in the required directions. No numerical or
symbolic computation is used.

A search by the exact source title and open-problem phrase, and by Grand
Lebesgue Sobolev/convolution sharpness keywords, did not locate these endpoint
lower bounds. The later interpolation paper arXiv:1705.06170 gives Young-type
upper estimates in related spaces but does not supply this lower-bound
argument. This should be regarded as a candidate new partial theorem pending
expert literature review.

The packet does not determine exact constants, treats the trace in the
Euclidean model geometry, and treats convolution on $\R^d$ rather than every
unimodular group. Arbitrary generating functions without regular endpoint
behavior can be governed by a different Legendre envelope. Human review is
recommended, especially for the source's intended projection convention and
for possible overlap with limiting interpolation results.

\begin{thebibliography}{9}
\bibitem{LOS}
E. Liflyand, E. Ostrovsky, and L. Sirota,
\emph{Tensor, Sobolev, Multiplicative and Convolution Operators in the
Bide-Side Grand Lebesque Spaces}, arXiv:0801.3404 (2008).
\bibitem{BL}
H. J. Brascamp and E. H. Lieb,
\emph{Best constants in Young's inequality, its converse, and its
generalization to more than three functions}, J. Funct. Anal. 20 (1976),
151--173.
\bibitem{FMDS}
P. Fern\'andez-Mart\'inez and E. Brandani da Silva,
\emph{New Young inequalities and applications}, arXiv:1705.06170;
Z. Anal. Anwend. 38 (2019), 419--437.
\end{thebibliography}
\end{document}
